% !TEX root = ../my-thesis.tex
%
\chapter{Conclusion}
\label{sec:conclusion}

This thesis presented an end-to-end pipeline that turns natural-language debate
text into a formally evaluated Dung argumentation framework, from categorized
attack extraction through confidence filtering to a Prolog solver under grounded,
preferred, and stable semantics.

Returning to the research questions, RQ1 is answered in the affirmative. The
pipeline runs end to end across eight topics and produces well-defined frameworks
whose extensions can be computed in a fraction of a second per topic thanks to
the grounded-first reduction. RQ4 is answered by the cross-topic analysis. All
frameworks are cyclic, and six of eight admit several preferred and stable
extensions, so the formal layer can make explicit several internally consistent
positions that the raw argument set leaves implicit, although its advantage over
the grounded extension is uneven across topics. The answers to RQ2 and RQ3 rest
on the typed-versus-binary ablation and the gold-standard evaluation, which are
set up in Chapter~\ref{sec:evaluation} and will be completed before submission.

% TODO: nach Durchfuehrung der Experimente die Antworten auf RQ2 und RQ3 hier
% konkret zusammenfassen.

Overall the results demonstrate the feasibility of integrating modern language
models with established formal argumentation, and they suggest that formal
semantics can provide a transparent reasoning layer on top of LLM-generated
argumentative content.

\section{Future Work}
\label{sec:conclusion:future}

Several directions remain open. The extraction quality should be evaluated
against annotated subtask datasets beyond the single gold standard used here.
The typed attacks and confidence scores, which the current pipeline discards
when converting to binary facts, could be used directly by a richer formalism
such as ASPIC+ or a quantitative bipolar framework. Finally, the pipeline should
be compared qualitatively against other methods, including a standalone LLM asked
to summarise the positions in a debate, and across several models rather than a
single one.
